\paragraph{Related Work}
The friction between \LaTeX{} and LLM-driven workflows is well-documented across three critical dimensions.

First, \textbf{generation and execution failures}. Benchmarks reveal that LLMs severely struggle with \LaTeX{}'s verbose syntax and long-range dependencies, showing low accuracy~\cite{kale2025texpert} and high token inefficiency in multimodal tasks~\cite{ling_table2latex_2025, xia2024docgenome}. This causes systemic bottlenecks in automated science pipelines~\cite{lu2024aiscientist} and ``execution illusions'' of uncompilable code~\cite{lyn2025translatex}. Thus, augmenting compiler logs with live ASTs is necessary for reliable AI assistance~\cite{jain2026bibby}.

Second, \textbf{training and representation overhead}. Upstream, processing heterogeneous \LaTeX{} incurs massive engineering overhead~\cite{paster2023openwebmath} and prohibitive token costs~\cite{lin2024accurate}. \LaTeX{}'s syntactic redundancy reduces predictability and inflates training costs~\cite{taylor2022galactica}. Conversely, explicitly modeling markup hierarchies drastically improves document understanding~\cite{li2022markuplm} and AST-guided code generation~\cite{zeng2026treediffastguidedcodegeneration}.

Third, \textbf{system architecture limitations}. To address batch compilation delays, existing solutions propose incremental compilers~\cite{madje2023typst, haug2022typst} or GUI overlays~\cite{gobert2022ilatex}, but still separate authoring from visual output. Meanwhile, studies highlight \LaTeX{}'s real-world fragility~\cite{tan2024inconsistencies}, user inefficiency~\cite{knauff2014efficiency}, and extreme runtime costs when used as a computing substrate~\cite{gardner2025neuralatex}. We build on the tree-based WYSIWYG lineage of \TeX{}macs~\cite{van_der_hoeven_gnu_2001}, but our contribution is distinct: whereas \TeX{}macs is prior art designed for human editing, we treat its \texttt{.tmu} format as a low-entropy, tree-structured \emph{representation for LLMs} and contribute the first evaluation of such a structured representation on LLM generation and fine-tuning. In doing so, we provide evidence that a tree-structured format offers advantages over \LaTeX{} along both axes the literature above identifies as broken---structural transparency for generation and predictability for training. \textit{(A comprehensive review is deferred to Appendix~\ref{app:rew}.)}